\section{Shell脚本}
\subsection{脚本格式}
在 Shell 脚本 的第一行通常会写上\verb|#!/bin/bash|，
这行语句称为 shebang，告诉操作系统在执行脚本时，应该调用哪一个解释器。

\subsection{变量}
在 Shell 中，变量可以像 Python 一样直接赋值，不过这种方式定义的变量都是全局变量。
除了直接赋值，还可以通过 read 从标准输入读取用户输入并赋值，一般配合 echo -n 给出提示，
也可以通过\verb|$0、$1|等读取命令行参数，\verb|$#|表示参数的个数，\verb|$@|表示参数列表。

\begin{code}[title=变量]
	\begin{verbatim}
local foo="a"       # 局部变量
echo $REPLY         # read 未指定变量时的默认存储变量
declare -u upper    # 声明变量，赋值时自动转大写
declare -l lower    # 声明变量，赋值时自动转小写
read -p "prompt: " -t 5 var         # 限时5s
read -p "password: " -s var           # -s隐藏输入
read -e -i "default" var          # -e命令补全 -i提供默认值
	\end{verbatim}
\end{code}

在需要向命令传递输入时，Shell 还提供了 两种重定向方式：here string和here document。
Here String 使用 \verb|<<<|，可以把一行字符串作为标准输入传递给命令；
Here Document则适合多行输入，它把分隔符之间的所有文本传递给命令的标准输入。
\begin{code}[title=Here Document]
	\begin{verbatim}
# 其中 EOF 是分隔符，可以自定义，只要首尾一致即可。
cat <<EOF
这是第一行
这是第二行
EOF
	\end{verbatim}
\end{code}

\subsection{语句}
\begin{code}[title=语句]
	\begin{verbatim}
[ -e file ]              # 文件存在（-d目录 -f普通文件 -r可读  -s非空）
[[ str1 >= str2 ]]       # 字符串字母序比较，== 支持 glob 通配符（* ? [...]）
[[ str =~ ^[0-9]+$ ]]   # =~ 支持正则表达式
(( INT == 0 ))           # 整数比较

for (( i=0; i<6; i=i+1 )); do
    echo $i
done

$((2#1010))          # 2 进制
$((8#77))            # 8 进制（0 前缀亦可）
$((16#ff))           # 16 进制（0x 前缀亦可）
$((2**10))           # 求幂，其余运算符与 C 语言相同
	\end{verbatim}
\end{code}

\subsection{参数展开}
\texttt{\#} 表示从头匹配，\texttt{\%} 表示从尾匹配；双符号 \texttt{\#\#} / \texttt{\%\%} 为最长匹配，单符号为最短匹配。
\begin{code}[title=参数展开]
	\begin{verbatim}
"${a}_file"          # 变量与字符串拼接
${a:-"test"}         # a 未设置时用 test 替代
${a:+"test"}         # a 已设置时用 test 替代
${a:="test"}         # a 未设置时用 test 替代并同时赋值给 a
${a:?"error"}        # a 未设置时返回错误

${#a}                # 字符串长度
${a:offset}          # 从 offset 开始的子串
${a:offset:length}   # 从 offset 开始长度为 length 的子串
${a,,}               # 转小写
${a^^}               # 转大写
${a^}                # 首字母大写

${a#pattern}         # 从头删除最短匹配
${a%%pattern}        # 从尾删除最长匹配
${a/pattern/string}  # 替换第一个匹配
${a//pattern/string} # 替换所有匹配
${a/#pattern/string} # 替换开头匹配
${a/%pattern/string} # 替换结尾匹配
	\end{verbatim}
\end{code}
